\clearpage\SourceLeaf{111}%
\pdfbookmark[1]{Über die Grenzen der geozentrischen Positionen der Planeten}{grenzen1804}
\begin{center}
{\small ÜBER DIE GRENZEN DER GEOZENTRISCHEN POSITIONEN DER PLANETEN.\par}
\medskip\hrule\medskip
{\small Monatliche Correspondenz zur Beförderung der Erd- und Himmels-Kunde,\par
herausgegeben vom Freiherrn \textsc{von Zach}. August 1804.\par}
\medskip\hrule
\end{center}\medskip
Von der Sonne aus gesehen scheint die Bewegung jedes Planeten immer
auf ein und demselben Großkreis am Fixsternhimmel zu verlaufen, wenn man
von den kleinen Störungen durch andere Himmelskörper absieht. Von der
Erde aus würde sie ebenso erscheinen, wenn die Ebene seiner Bahn mit
der Ekliptik zusammenfiele. Sind diese beiden Ebenen jedoch gegeneinander
geneigt, so liegen alle möglichen geozentrischen Positionen des Planeten
auf der Himmelskugel nicht mehr wie im ersten Fall auf einer nur in
einer Dimension ausgedehnten Linie. Sie füllen vielmehr einen
Flächenbereich, eine \textit{Zone}, die den ganzen übrigen Himmel, an
dem der Planet von der Erde aus nie erscheinen kann, in zwei Teile
trennt und daher der \textit{Zodiakus} des Planeten heißen kann. Somit
hat jeder Planet im Grunde seinen eigenen Zodiakus, dessen Grenzen
(\textit{Limiten}) sich vollkommen genau bestimmen lassen, wenn man
seine Bahn und die Erdbahn als Kegelschnitte mit unveränderlichen
Elementen auffasst. Schon für sich genommen ist die genaue Bestimmung
dieser Grenzen ein interessantes analytisches Problem. Auch ihre
Anwendung hat einige praktische Bedeutung, besonders bei den beiden neuen
Planeten, deren Zodiaken eine beträchtliche Ausdehnung haben. Bekanntlich
braucht man sehr genaue und detaillierte Sternkarten, um diese
bemerkenswerten Himmelskörper aufzusuchen und zu beobachten. Selbst die
besten Karten, die wir bisher
\SourceLeaf{112}%
besitzen, reichen dafür noch bei Weitem nicht aus. Will man also nicht
jedes Jahr Spezialkarten der Gegenden entwerfen, die diese Planeten
durchlaufen, so muss man einen eigenen Atlas in Betracht ziehen, der
den ganzen Bereich umfasst, in dem sie erscheinen können. Die genaue
Bestimmung seiner Grenzen ist umso wünschenswerter, als man sich bei
einem Vorhaben von ohnehin beträchtlichem Umfang gerne jede dafür
unnötige Mühe erspart. Gewiss wird allen Freunden der Astronomie die
Nachricht sehr willkommen sein, dass der geschickte Lilienthaler
Astronom \textsc{Harding}, dem wir bereits mehrere vortreffliche
Spezialkarten verdanken, mit dieser größeren Arbeit begonnen hat. Sie
wird sich sowohl durch ihren großen Detailreichtum als auch durch die
besonders sorgfältige, durchweg auf eigener Anschauung beruhende
kritische Prüfung auszeichnen.

Wir legen durch den Mittelpunkt der Sonne drei zueinander senkrechte,
im Übrigen beliebige Ebenen. Die senkrechten Abstände des beobachteten
Planeten von ihnen nennen wir $x$, $y$, $z$, die Abstände der Erde dagegen
$x'$, $y'$, $z'$. Ferner setzen wir
\[
\begin{aligned}
x'-x&=\Delta\cos\alpha\cos\ObservedBeta\\
y'-y&=\Delta\sin\alpha\cos\ObservedBeta\\
z'-z&=\Delta\sin\ObservedBeta
\end{aligned}
\]
Dabei ist $\Delta$ der Abstand des Planeten von der Erde,
$\ObservedBeta$ die Neigung der vom Planeten zur Erde gezogenen Geraden
gegen eine durch die Erde gelegte Ebene parallel zur Ebene der
$z$-Koordinaten und $\alpha$ der Winkel zwischen der Projektion dieser
Geraden auf diese Ebene und einer durch die Erde gelegten Ebene parallel
zur Ebene der $y$-Koordinaten. Auf der Himmelskugel bestimmen also $\alpha$
und $\ObservedBeta$ die geozentrische Position des Planeten in Bezug auf die
Großkreise, die von den Ebenen der $z$-, $y$-, $x$-Koordinaten (genauer:
von den dazu parallelen Ebenen durch die Erde) gebildet werden. Ebenso
bestimmen Länge und Breite die Position in Bezug auf die Ekliptik und
die Äquinoktial- und Solstitialkoluren.

Aus der Beschaffenheit der Planetenbahn ergeben sich zwei Gleichungen
zwischen $x$, $y$, $z$. Daher kann man diese drei veränderlichen Größen
als Funktionen \textit{einer} veränderlichen Größe betrachten, die wir
mit $t$ bezeichnen und im Übrigen noch offenlassen. Ebenso sollen $x'$,
$y'$, $z'$ Funktionen der veränderlichen Größe
\SourceLeaf{113}%
$t'$ sein. Damit sind $\alpha$ und $\ObservedBeta$ Funktionen beider
Größen $t$, $t'$. Sie mögen durch die folgenden Differentialgleichungen
bestimmt sein.
\[
\mathrm{d}\alpha=p\mathrm{d}t+p'\mathrm{d}t',\qquad \mathrm{d}\ObservedBeta=q\mathrm{d}t+q'\mathrm{d}t'
\]

Unter dieser Voraussetzung ist klar, dass bei gleichzeitiger Änderung
von $t$, $t'$ im Verhältnis $\mathrm{d}t:\mathrm{d}t'=-p':p$ die Größe
$\alpha$ unverändert bleibt, während $\ObservedBeta$ so lange zu- oder
abnimmt, bis es einen größten oder kleinsten Wert erreicht hat. Das
geschieht offenbar, wenn $pq'-qp'=0$ wird. Die Verbindung \textit{aller}
Werte von $t$ und $t'$ ergibt nun alle möglichen geozentrischen
Positionen des Planeten. Unter allen Verbindungen, die dasselbe
$\alpha$ ergeben, muss diejenige eintreten, bei der $\ObservedBeta$
einen größten oder kleinsten Wert annimmt, wenn die geozentrische
Position auf den Grenzen des Zodiakus des Planeten liegen soll.
Daraus folgt, dass die Bedingungsgleichung $pq'-qp'=0$ diese Grenzen
bestimmt.

Durch Differenzieren der obigen Gleichungen erhält man
\[
\begin{aligned}
\mathrm{d}x'-\mathrm{d}x&=\cos\alpha\cos\ObservedBeta\,\mathrm{d}\Delta-\Delta\sin\alpha\cos\ObservedBeta\,\mathrm{d}\alpha-\Delta\cos\alpha\sin\ObservedBeta\,\mathrm{d}\ObservedBeta\\
\mathrm{d}y'-\mathrm{d}y&=\sin\alpha\cos\ObservedBeta\,\mathrm{d}\Delta+\Delta\cos\alpha\cos\ObservedBeta\,\mathrm{d}\alpha-\Delta\sin\alpha\sin\ObservedBeta\,\mathrm{d}\ObservedBeta\\
\mathrm{d}z'-\mathrm{d}z&=\sin\ObservedBeta\,\mathrm{d}\Delta+\Delta\cos\ObservedBeta\,\mathrm{d}\ObservedBeta
\end{aligned}
\]

Zusammen mit jenen Gleichungen folgt daraus leicht
\[
\begin{gathered}
-\sin\alpha(\mathrm{d}x'-\mathrm{d}x)+\cos\alpha(\mathrm{d}y'-\mathrm{d}y)=\Delta\cos\ObservedBeta\,\mathrm{d}\alpha\\
\begin{aligned}
-\cos\alpha\sin\ObservedBeta(\mathrm{d}x'-\mathrm{d}x)-\sin\alpha\sin\ObservedBeta(\mathrm{d}y'-\mathrm{d}y)&\\
{}+\cos\ObservedBeta(\mathrm{d}z'-\mathrm{d}z)&=\Delta\mathrm{d}\ObservedBeta
\end{aligned}
\end{gathered}
\]

Aus den partiellen Differentialen ergibt sich also
\[
\begin{aligned}
\Delta\cos\ObservedBeta\,p\mathrm{d}t&=\sin\alpha\,\mathrm{d}x-\cos\alpha\,\mathrm{d}y\\
\Delta\cos\ObservedBeta\,p'\mathrm{d}t'&=-\sin\alpha\,\mathrm{d}x'+\cos\alpha\,\mathrm{d}y'\\
\Delta q\mathrm{d}t&=\cos\alpha\sin\ObservedBeta\,\mathrm{d}x+\sin\alpha\sin\ObservedBeta\,\mathrm{d}y-\cos\ObservedBeta\,\mathrm{d}z\\
\Delta q'\mathrm{d}t'&=-\cos\alpha\sin\ObservedBeta\,\mathrm{d}x'-\sin\alpha\sin\ObservedBeta\,\mathrm{d}y'+\cos\ObservedBeta\,\mathrm{d}z'
\end{aligned}
\]

Setzt man diese Werte von $p$, $p'$, $q$, $q'$ in die Bedingungsgleichung
$pq'=p'q$ ein, so erhält man nach den nötigen Umformungen
\[
\begin{aligned}
\cos\alpha\cos\ObservedBeta(\mathrm{d}y'\mathrm{d}z-\mathrm{d}y\,\mathrm{d}z')&+\sin\alpha\cos\ObservedBeta(\mathrm{d}z'\mathrm{d}x-\mathrm{d}z\,\mathrm{d}x')\\
&+\sin\ObservedBeta(\mathrm{d}x'\mathrm{d}y-\mathrm{d}x\,\mathrm{d}y')=0
\end{aligned}
\]
\SourceLeaf{114}%
oder nach Multiplikation mit $\Delta$
\[
\begin{aligned}
(x'-x)(\mathrm{d}y'\mathrm{d}z-\mathrm{d}y\,\mathrm{d}z')&+(y'-y)(\mathrm{d}z'\mathrm{d}x-\mathrm{d}z\,\mathrm{d}x')\\
&+(z'-z)(\mathrm{d}x'\mathrm{d}y-\mathrm{d}x\,\mathrm{d}y')=0
\end{aligned}
\]
Diese Gleichung lässt sich noch besser in folgender Form schreiben.
\[
\left.\begin{aligned}
\mathrm{d}x(y'\mathrm{d}z'-z'\mathrm{d}y')&+\mathrm{d}x'(y\,\mathrm{d}z-z\,\mathrm{d}y)\\
{}+\mathrm{d}y(z'\mathrm{d}x'-x'\mathrm{d}z')&+\mathrm{d}y'(z\,\mathrm{d}x-x\,\mathrm{d}z)\\
{}+\mathrm{d}z(x'\mathrm{d}y'-y'\mathrm{d}x')&+\mathrm{d}z'(x\,\mathrm{d}y-y\,\mathrm{d}x)
\end{aligned}\right\}=0
\]

Diese Gleichung gibt allgemein den Zusammenhang zwischen den
Positionen der Erde und des Planeten an, bei denen die geozentrische
Position des Planeten auf den Grenzen liegt. Man muss darin nur für
$x$, $y$, $z$ ihre Werte in Abhängigkeit von $t$ und für $x'$, $y'$, $z'$
ihre Werte in Abhängigkeit von $t'$ einsetzen, wie sie sich jeweils
aus der Bahn ergeben, um eine endliche Gleichung zwischen $t$ und $t'$
zu erhalten. Es schien die Mühe wert, diese Gleichung durch eine
allgemeine Analyse herzuleiten. Es ist aber leicht zu zeigen, dass sie
zugleich die Bedingung dafür ist, \textit{dass die Tangenten an den
Positionen der Erde und des Planeten in einer Ebene liegen}, und gerade
diese Bedingung aus den Anforderungen unserer Aufgabe abzuleiten.
Der Kürze halber gehen wir hier jedoch nicht weiter darauf ein.

Für die Größen $t$, $t'$, die wir bisher offengelassen haben, wählen wir
am bequemsten die heliozentrischen Winkelabstände des Planeten und der
Erde auf ihren Bahnen von der gemeinsamen Knotenlinie, und zwar vom
aufsteigenden Knoten der Planetenbahn auf der Erdbahn. Bezeichnen wir
nun die Abstände des Planeten und der Erde von der Sonne mit $r$, $r'$,
so lassen sich die Koordinaten wie folgt ausdrücken.
\[
\begin{aligned}
x&=r\sin a\sin(t+A)\\
y&=r\sin b\sin(t+B)\\
z&=r\sin c\sin(t+C)\\
x'&=r'\sin a'\sin(t'+A')\\
y'&=r'\sin b'\sin(t'+B')\\
z'&=r'\sin c'\sin(t'+C')
\end{aligned}
\]

Dies und die Bedeutung der Konstanten $a$, $A$ usw. lassen sich leicht
durch Verallgemeinerung der in der \textit{Mon. Corr.} vom Mai 1804
[S.\,98 d.\,B.] vorgetragenen Untersuchung verstehen. Sind ferner $k$, $k'$
\SourceLeaf{115}%
die Halbparameter der Kegelschnitte, die der Planet und die Erde
beschreiben, $e$, $e'$ die Exzentrizitäten und $g$, $g'$ die
Winkelabstände der Aphele von der Knotenlinie, so ergibt sich
\[
r=\frac{k}{1-e\cos(t-g)},\qquad r'=\frac{k'}{1-e'\cos(t'-g')}
\]

Die Rechnung ergibt daraus
\[
\mathrm{d}x=\frac{k\sin a\,\mathrm{d}t}{(1-e\cos(t-g))^2}\times(\cos(t+A)-e\cos(g+A))
\]

Die Werte von $\mathrm{d}y$, $\mathrm{d}z$ haben eine ähnliche Form.
Um sie zu erhalten, braucht man nur $a$, $A$ durch $b$, $B$ oder durch
$c$, $C$ zu ersetzen. Die Werte von $\mathrm{d}x'$, $\mathrm{d}y'$,
$\mathrm{d}z'$ erhält man aus denen von $\mathrm{d}x$, $\mathrm{d}y$,
$\mathrm{d}z$, indem man die auf den Planeten bezogenen Größen durch
die entsprechenden Größen für die Erde ersetzt.

Bequemer entwickelt man $y\,\mathrm{d}z-z\,\mathrm{d}y$ aus den
Werten von $y$, $z$, bevor man darin den Wert von $r$ einsetzt. So
erhält man
\[
y\,\mathrm{d}z-z\,\mathrm{d}y=rr\sin b\sin c\sin(B-C)\mathrm{d}t=rr\cos a\,\mathrm{d}t
\]
(siehe den genannten Aufsatz [S.\,100 d.\,B.]), und ebenso
\[
\begin{aligned}
z\,\mathrm{d}x-x\,\mathrm{d}z&=rr\cos b\,\mathrm{d}t\\
x\,\mathrm{d}y-y\,\mathrm{d}x&=rr\cos c\,\mathrm{d}t
\end{aligned}
\]

Für die drei entsprechenden Ausdrücke für die Erde erhält man
ganz ähnliche Werte.

Setzt man alle diese Werte ein und formt entsprechend um, so
erhält die obige Bedingungsgleichung die folgende Gestalt.
\[
\begin{aligned}
&k'\cos a'\sin a\,(\cos(t+A)-e\cos(g+A))\\
&{}+k'\cos b'\sin b\,(\cos(t+B)-e\cos(g+B))\\
&{}+k'\cos c'\sin c\,(\cos(t+C)-e\cos(g+C))\\
&{}+k\cos a\sin a'\,(\cos(t'+A')-e'\cos(g'+A'))\\
&{}+k\cos b\sin b'\,(\cos(t'+B')-e'\cos(g'+B'))\\
&{}+k\cos c\sin c'\,(\cos(t'+C')-e'\cos(g'+C'))=0
\end{aligned}
\]

Bezeichnet man die Neigung der Planetenbahn gegen die Erdbahn mit
$i$, so lassen sich die ersten drei Terme dieser Gleichung durch
geeignete Umformungen in $k'\sin i\,(\cos t-e\cos g)$ und die letzten
drei in $-k\sin i\,(\cos t'-e'\cos g')$ verwandeln.
\SourceLeaf{116}%
Die Umformungen sind an sich leicht, wenn auch etwas umständlich.
Wir können uns hier ihre Ausführung umso eher ersparen, als wir zum
selben Ergebnis viel bequemer gelangen, wenn wir die drei bisher
offengelassenen Grundebenen, auf die sich die Koordinaten beziehen,
geeignet wählen. Wir nehmen nämlich für die Ebene der $z$-Koordinaten
die Ekliptik und wählen die Ebenen der $x$-, $y$-Koordinaten so, dass
der Pol der ersten in den aufsteigenden Knoten der Planetenbahn fällt,
der Pol der zweiten dagegen $90^0$ weiter vorwärts auf der Ekliptik
liegt. Auf der Seite \textit{dieser} Pole und auf der Nordseite der
Ekliptik sollen die Koordinaten $x$, $y$, $z$ positiv angesetzt werden.
Man sieht leicht, dass unter diesen Voraussetzungen gilt
\[
\begin{array}{llll}
a=90^0,&b=90^0+i,&c=i,&C=0\\
a'=90^0,&b'=90^0.&c'=0,&B'=0
\end{array}
\]
und somit die obige Gleichung übergeht in
\[
k'\sin i\,(\cos t-e\cos g)-k\sin i\,(\cos t'-e'\cos g')=0
\]
oder
\[
k'(\cos t-e\cos g)=k(\cos t'-e'\cos g')
\]

Aus der Theorie der Kegelschnitte lässt sich leicht zeigen, dass
$\frac{k}{\cos t-e\cos g}$ und $\frac{k'}{\cos t'-e'\cos g'}$ die
Abstände der Sonne von den Schnittpunkten der Knotenlinie der beiden
Bahnen mit den Tangenten an den Positionen des Planeten und der Erde
angeben. Die soeben gefundene Gleichung zeigt also, dass diese beiden
Tangenten die Knotenlinie in ein und demselben Punkt schneiden. Das
stimmt mit der oben erwähnten Bedingung überein, nach der sie in ein
und derselben Ebene liegen sollen.

Bei der Lage der Planetenbahn zur Erdbahn sind drei Fälle zu
unterscheiden. Entweder schließt die Planetenbahn die Erdbahn ein,
oder die Erdbahn die Planetenbahn, oder beide umschließen sich wie
Glieder einer Kette. Der erste Fall tritt ein, wenn der Planet auf
der Knotenlinie auf beiden Seiten weiter von der Sonne entfernt ist
als die Erde. Der zweite tritt ein, wenn die Erde auf beiden Seiten
weiter von der Sonne entfernt ist als der Planet. Der dritte tritt
ein, wenn auf der einen Seite der Planet und auf der anderen die
Erde weiter von der Sonne entfernt ist. Keiner der bisher bekannten
Planeten hat gegenüber der Erde oder einem anderen Planeten die Lage,
die der dritte Fall erfordert.
\SourceLeaf{117}%
Kometen dieser Art gibt es dagegen in großer Zahl. Die analytische
Bedingung für den ersten Fall ist, dass $k-k'$ positiv und ohne
Rücksicht auf das Vorzeichen größer ist als $k'e\cos g-ke'\cos g'$.
Im zweiten Fall muss $k'-k$ diese Eigenschaften haben. Im dritten
Fall muss $k-k'$ oder $k'-k$ ohne Rücksicht auf das Vorzeichen kleiner
sein als $k'e\cos g-ke'\cos g'$.

Im ersten dieser drei Fälle erhält man aus der obigen Gleichung zu
jedem beliebigen Wert von $t$ zwischen $0^0$ und $360^0$ zwei Werte
von $t'$. Im zweiten Fall liefert jeder Wert von $t'$ zwei Werte
von $t$. Der eine Wert von $t'$ (im ersten Fall, bzw. von $t$ im
zweiten) liegt stets zwischen $0$ und $180$, der andere zwischen $180$
und $360^0$, oder vielmehr innerhalb noch engerer Grenzen, die sich
leicht genauer bestimmen lassen. Im ersten Fall entsprechen also
jeder heliozentrischen Position des Planeten zwei heliozentrische
Positionen der Erde. Umgekehrt gilt das nicht, denn die Erde kann den
Planeten nur in zwei getrennten Abschnitten ihrer Bahn auf seinen
Grenzen sehen. Der eine liegt unterhalb oder südlich der Ebene der
Planetenbahn, der andere oberhalb oder nördlich davon. Auf der
nördlichen Grenze ist der Planet nur zu sehen, wenn $t'$ einen seiner
möglichen Werte zwischen $0$ und $180$ annimmt, auf der südlichen bei
einem Wert zwischen $180$ und $360^0$. Ebenso entsprechen im zweiten
Fall jeder heliozentrischen Position der Erde zwei heliozentrische
Positionen des Planeten. Umgekehrt gilt das auch hier nicht, denn der
Planet kann der Erde nur in zwei getrennten Abschnitten seiner Bahn
auf seinen Grenzen erscheinen, von denen der eine nördlich, der
andere südlich der Ekliptik liegt. Auf der nördlichen Grenze erscheint
er bei den Werten zwischen $0$ und $180$, auf der südlichen bei den
Werten zwischen $180$ und $360^0$.

Im dritten Fall dagegen können $t$ und $t'$ jeweils nur Werte
innerhalb bestimmter Grenzen annehmen. Anders gesagt, Erde und
Planet müssen sich jeweils in einem bestimmten Abschnitt ihrer
Bahn befinden, wenn die obige Bedingungsgleichung erfüllt sein soll.

Daraus ergibt sich, dass die geozentrischen Positionen des Planeten,
die aus allen der obigen Gleichung genügenden Verbindungen seiner
heliozentrischen Positionen mit denen der Erde hervorgehen, im ersten
und zweiten Fall zwei getrennte, jeweils in sich zurücklaufende Linien
auf der Himmelskugel bilden. Zwischen ihnen liegt im ersten Fall der
Großkreis, der die Ebene der Planetenbahn darstellt, im zweiten die
Ekliptik.
\SourceLeaf{118}%
Im dritten Fall hingegen bilden jene geozentrischen Positionen nur
\textit{eine} in sich zurücklaufende Linie, wie eine nähere Betrachtung
des Falls leicht zeigt.

Nach den vorhergehenden Untersuchungen kann der Zodiakus des Planeten
keine anderen Grenzen haben als eben diese Linien. Es scheint daher
naheliegend, dass der Zodiakus des Planeten in den ersten beiden
Fällen die zwischen den beiden Linien liegende Zone ist, im dritten
Fall einer der beiden Bereiche, in die jene Linie die ganze
Kugeloberfläche teilt. Dieser Schluss träfe jedoch in den ersten
beiden Fällen nicht immer und im dritten nie zu. Man darf hier,
wie in vielen anderen Anwendungen der Analysis, bei denen eine
ähnliche Vorsicht mitunter zu wenig beachtet wird, nicht vergessen,
dass unsere Schlussfolgerung ganz auf der Voraussetzung beruht, der
Zodiakus des Planeten sei tatsächlich begrenzt und der Planet könne
von der Erde aus nicht an jedem Punkt des Himmels erscheinen. Diese
Voraussetzung ist jedoch im dritten Fall nicht erfüllt, wie sich
schon aus Gründen der \textit{Geometrie der Lage} zeigen lässt. Die gefundene
Linie kann hier also keine Grenze des Zodiakus des Planeten sein,
da dieser den \textit{ganzen} Himmel einnimmt. Im ersten und zweiten
Fall gibt es zwar immer wenigstens auf einer Seite der gefundenen
Zone Stellen am Himmel, an denen der Planet nie erscheinen kann.
Damit ist die eine Linie gewiss eine Grenze. Trotzdem kann es
vorkommen, dass es \textit{nur} auf einer Seite solche ausgeschlossenen
Stellen gibt. Dann bildet die andere Linie keine Grenze, und der
von ihr von der Zone abgetrennte Bereich des Himmels gehört ebenso
ganz zum Zodiakus wie die Zone selbst. Hier ist jedoch nicht der Ort,
diese Untersuchung vollständig durchzuführen und darzulegen, was
jene Linien in solchen Fällen eigentlich bedeuten, wenn sie keine
Grenzen sind. Wir können uns hier umso eher damit begnügen, die
Freunde der Analysis auf diese scheinbar paradoxen Erscheinungen
aufmerksam zu machen, als sich leicht zeigen lässt, dass keiner der
bisher bekannten Planeten, auf die es uns hier zunächst ankommt,
von der Erde aus je an einem der beiden Pole der Ekliptik erscheint.
Die genannten Ausnahmen können bei diesen Planeten also nicht
eintreten. Ihre Zodiaken müssen somit wirkliche \textit{Zonen} sein
und die beiden gefundenen Linien ihre Grenzen.

Wir wollen unsere Ergebnisse nun auf \textit{Pallas} und \textit{Ceres}
anwenden, teils zur weiteren Erläuterung, teils für den praktischen
Gebrauch.
\SourceLeaf{119}%
Wir wollen die Grenzen ihrer Zodiaken so abstecken, dass man sie in
die Sternkarten eintragen kann. Wegen der Störungen müssen diese
Grenzen noch etwas erweitert werden, und weil sich die Elemente in
Zukunft ändern werden, sind einige Änderungen nötig. Für die Praxis
werden sie jedoch unerheblich sein. Wenn die Beobachtungen sie nach
Jahren bemerkbar machen, werden sie gerade deshalb sogar interessant
sein, weil sie dann deutlich zeigen, wie die hier zugrunde gelegten
elliptischen Elemente allmählich nicht mehr ausreichen.

Für \textit{Pallas} setzen wir nach den neuesten Elementen für den
Anfang des Jahres 1803 (\textit{Mon. Corr.}, März 1804)
\[
\begin{array}{rl@{\qquad}rl}
e&=0.2457396&\varphi&=14^0\;13'\;31''966\\
k&=2.602122&&\\
g&=128^0\;49'\;20''7&&\\
\AscendingNode&=172\;28\;13.7&i&=34\;38\;1.1
\end{array}
\]

Für die Erde dagegen
\[
\begin{array}{rl@{\qquad}rl}
e'&=0.016792&\varphi&=0^0\;57'\;43''802\\
k'&=0.999718&&\\
g'&=107^0\;5'\;39''&&\\
\mathrm{Aph.}&=279\;33\;52.7&&
\end{array}
\]

Setzt man diese Werte ein, so erhält unsere obige Gleichung die Form
\[
\cos t'=0.384193\cos t+0.0542514
\]

Daraus folgt, dass die beiden äußersten Werte von $\cos t'$
$+0.438444$ und $-0.329942$ sind. Alle möglichen Werte von $t'$
liegen also entweder zwischen $63^0\;59'\;43''$ und $109^0\;15'\;55''$
oder zwischen $250^0\;44'\;5''$ und $296^0\;0'\;17''$. Daher kann
\textit{Pallas} der Erde nur dann auf ihren Grenzen erscheinen,
wenn die heliozentrische Länge der Erde zwischen $236^0\;28'$
und $281^0\;44'$ oder zwischen $63^0\;12'$ und $180^0\;29'$ liegt,
also etwa vom 18.\,Mai bis 4.\,Juli und vom 26.\,November bis
9.\,Januar. Im ersten Abschnitt ihrer Bahn liegt die Erde südlich,
im anderen nördlich der Ebene der Pallasbahn. Daher wird ihr
\textit{Pallas} im ersten Abschnitt auf der nördlichen, im anderen
auf der südlichen Grenze ihres Zodiakus erscheinen. Es lässt sich
auch leicht zeigen, dass \textit{Pallas}
\SourceLeaf{120}%
jedes Jahr zu den angegebenen Zeiten je einmal die nördliche und
die südliche Grenze berühren muss. --- Um genügend Punkte auf beiden
Grenzen zu erhalten, nehmen wir für $t$ der Reihe nach alle Werte
von $0$ bis $360^0$ von $10$ zu $10$ Grad. Kombinieren wir
jeden davon mit den beiden zugehörigen, aus der obigen Formel
zu bestimmenden Werten von $t'$, so leiten wir unmittelbar die
entsprechenden geozentrischen Positionen in Rektaszension und
Deklination ab. Dazu können wir das im genannten Aufsatz erklärte
Verfahren und die dort bereits berechneten Konstanten verwenden.
Die folgende Tabelle enthält die Ergebnisse dieser Rechnungen.

\begin{center}
\begingroup\small\setlength{\tabcolsep}{5pt}
\begin{tabular}{r||rr|rrl||rr|rrl}
&\multicolumn{5}{c||}{Nördliche Grenze}&\multicolumn{5}{c}{Südliche Grenze}\\
$t$&\multicolumn{2}{c|}{Rektasz.}&\multicolumn{3}{c||}{Deklination}&\multicolumn{2}{c|}{Rektasz.}&\multicolumn{3}{c}{Deklination.}\\\hline
$0^0$&$148^0$&$4'$&$12^0$&$56'$&nördl.&$196^0$&$54'$&$7^0$&$12'$&südl.\\
$10$&$159$&$30$&$15$&$56$&&$205$&$54$&$4$&$26$&\\
$20$&$171$&$54$&$18$&$55$&&$214$&$15$&$2$&$11$&\\
$30$&$185$&$6$&$21$&$38$&&$222$&$3$&$0$&$23$&\\
$40$&$198$&$48$&$23$&$52$&&$229$&$27$&$1$&$1$&nördl.\\
$50$&$212$&$41$&$25$&$26$&&$236$&$32$&$2$&$4$&\\
$60$&$226$&$28$&$26$&$16$&&$243$&$26$&$2$&$48$&\\[3pt]
$70$&$239$&$59$&$26$&$23$&&$250$&$16$&$3$&$16$&\\
$80$&$253$&$8$&$25$&$53$&&$257$&$6$&$3$&$28$&\\
$90$&$265$&$54$&$24$&$49$&&$264$&$1$&$3$&$25$&\\
$100$&$278$&$18$&$23$&$19$&&$271$&$4$&$3$&$5$&\\
$110$&$290$&$25$&$21$&$27$&&$278$&$18$&$2$&$30$&\\
$120$&$302$&$19$&$19$&$17$&&$285$&$44$&$1$&$37$&\\[3pt]
$130$&$314$&$4$&$16$&$53$&&$293$&$23$&$0$&$27$&\\
$140$&$325$&$43$&$14$&$22$&&$301$&$16$&$1$&$2$&südl.\\
$150$&$337$&$16$&$11$&$47$&&$309$&$24$&$2$&$50$&\\
$160$&$348$&$42$&$9$&$14$&&$317$&$49$&$4$&$57$&\\
$170$&$359$&$55$&$6$&$49$&&$326$&$34$&$7$&$23$&\\
$180$&$10$&$49$&$4$&$40$&&$335$&$41$&$10$&$8$&\\[3pt]
$190$&$21$&$15$&$2$&$49$&&$345$&$12$&$13$&$10$&\\
$200$&$31$&$5$&$1$&$21$&&$355$&$12$&$16$&$24$&\\
$210$&$40$&$14$&$0$&$14$&&$5$&$44$&$19$&$53$&\\
$220$&$48$&$39$&$0$&$31$&südl.&$16$&$53$&$23$&$6$&\\
$230$&$56$&$23$&$0$&$59$&&$28$&$41$&$26$&$16$&\\
$240$&$63$&$28$&$1$&$10$&&$41$&$12$&$29$&$5$&\\[3pt]
$250$&$70$&$2$&$1$&$9$&&$54$&$25$&$31$&$21$&\\
$260$&$76$&$12$&$0$&$57$&&$68$&$16$&$32$&$52$&\\
$270$&$82$&$5$&$0$&$35$&&$82$&$36$&$33$&$29$&\\
$280$&$87$&$51$&$0$&$4$&&$97$&$12$&$33$&$5$&\\
$290$&$93$&$38$&$0$&$37$&nördl.&$111$&$47$&$31$&$39$&\\
$300$&$99$&$35$&$1$&$30$&&$126$&$6$&$29$&$14$&\\[3pt]
$310$&$105$&$53$&$2$&$36$&&$139$&$52$&$26$&$0$&\\
$320$&$112$&$40$&$3$&$59$&&$152$&$57$&$22$&$12$&\\
$330$&$120$&$8$&$5$&$40$&&$165$&$12$&$18$&$9$&\\
$340$&$128$&$26$&$7$&$44$&&$176$&$37$&$14$&$10$&\\
$350$&$137$&$43$&$10$&$10$&&$187$&$10$&$10$&$28$&\\
$360$&$148$&$4$&$12$&$56$&&$196$&$54$&$7$&$12$&
\end{tabular}
\endgroup
\end{center}
\SourceLeaf{121}%
Der Bequemlichkeit halber wurden mit geeigneten Interpolationsverfahren
zwischen diesen $72$ Punkten die folgenden $144$ Punkte eingefügt. Ihre
Rektaszensionen steigen von $5$ zu $5$ Grad.

\begin{center}
\begingroup\small\setlength{\tabcolsep}{3pt}
\begin{tabular}{r|rrl|rrl||r|rrl|rrl}
\multicolumn{14}{c}{\textit{Zodiakus der Pallas}}\\[4pt]
&\multicolumn{6}{c||}{Deklination}&&\multicolumn{6}{c}{Deklination}\\
\shortstack{Rektasz.}&\multicolumn{3}{c|}{\shortstack{der nördlichen\\Grenze}}&\multicolumn{3}{c||}{\shortstack{der südlichen\\Grenze}}&\shortstack{Rektasz.}&\multicolumn{3}{c|}{\shortstack{der nördlichen\\Grenze}}&\multicolumn{3}{c}{\shortstack{der südlichen\\Grenze}}\\\hline
$0^0$&$6^0$&$48'$&nördl.&$17^0$&$56'$&südl.&$180^0$&$20^0$&$38'$&nördl.&$12^0$&$58'$&südl.\\
$5$&$5$&$47$&&$19$&$32$&&$185$&$21$&$37$&&$11$&$12$&\\
$10$&$4$&$49$&&$21$&$4$&&$190$&$22$&$31$&&$9$&$29$&\\
$15$&$3$&$54$&&$22$&$33$&&$195$&$23$&$19$&&$7$&$49$&\\
$20$&$3$&$2$&&$23$&$59$&&$200$&$24$&$2$&&$6$&$13$&\\
$25$&$2$&$13$&&$25$&$20$&&$205$&$24$&$39$&&$4$&$42$&\\
$30$&$1$&$29$&&$26$&$36$&&$210$&$25$&$11$&&$3$&$17$&\\[3pt]
$35$&$0$&$50$&&$27$&$46$&&$215$&$25$&$37$&&$2$&$0$&\\
$40$&$0$&$16$&&$28$&$51$&&$220$&$25$&$58$&&$0$&$49$&\\
$45$&$0$&$14$&südl.&$29$&$49$&&$225$&$26$&$12$&&$0$&$13$&nördl.\\
$50$&$0$&$37$&&$30$&$41$&&$230$&$26$&$22$&&$1$&$7$&\\
$55$&$0$&$55$&&$31$&$26$&&$235$&$26$&$25$&&$1$&$52$&\\
$60$&$1$&$6$&&$32$&$4$&&$240$&$26$&$23$&&$2$&$28$&\\[3pt]
$65$&$1$&$11$&&$32$&$36$&&$245$&$26$&$16$&&$2$&$56$&\\
$70$&$1$&$9$&&$33$&$0$&&$250$&$26$&$3$&&$3$&$15$&\\
$75$&$1$&$0$&&$33$&$17$&&$255$&$25$&$45$&&$3$&$26$&\\
$80$&$0$&$44$&&$33$&$27$&&$260$&$25$&$23$&&$3$&$28$&\\
$85$&$0$&$21$&&$33$&$30$&&$265$&$24$&$55$&&$3$&$23$&\\
$90$&$0$&$10$&nördl.&$33$&$25$&&$270$&$24$&$22$&&$3$&$9$&\\[3pt]
$95$&$0$&$49$&&$33$&$13$&&$275$&$23$&$46$&&$2$&$48$&\\
$100$&$1$&$34$&&$32$&$54$&&$280$&$23$&$5$&&$2$&$19$&\\
$105$&$2$&$27$&&$32$&$27$&&$285$&$22$&$20$&&$1$&$43$&\\
$110$&$3$&$25$&&$31$&$53$&&$290$&$21$&$31$&&$1$&$0$&\\
$115$&$4$&$29$&&$31$&$11$&&$295$&$20$&$39$&&$0$&$10$&\\
$120$&$5$&$38$&&$30$&$23$&&$300$&$19$&$43$&&$0$&$46$&südl.\\[3pt]
$125$&$6$&$52$&&$29$&$27$&&$305$&$18$&$45$&&$1$&$49$&\\
$130$&$8$&$9$&&$28$&$24$&&$310$&$17$&$44$&&$2$&$58$&\\
$135$&$9$&$27$&&$27$&$14$&&$315$&$16$&$42$&&$4$&$12$&\\
$140$&$10$&$47$&&$25$&$57$&&$320$&$15$&$37$&&$5$&$32$&\\
$145$&$12$&$7$&&$24$&$35$&&$325$&$14$&$31$&&$6$&$56$&\\
$150$&$13$&$27$&&$23$&$7$&&$330$&$13$&$24$&&$8$&$24$&\\[3pt]
$155$&$14$&$46$&&$21$&$33$&&$335$&$12$&$17$&&$9$&$56$&\\
$160$&$16$&$4$&&$19$&$55$&&$340$&$11$&$10$&&$11$&$30$&\\
$165$&$17$&$19$&&$18$&$14$&&$345$&$10$&$3$&&$13$&$6$&\\
$170$&$18$&$29$&&$16$&$30$&&$350$&$8$&$57$&&$14$&$43$&\\
$175$&$19$&$36$&&$14$&$44$&&$355$&$7$&$52$&&$16$&$20$&\\
$180$&$20$&$38$&&$12$&$58$&&$360$&$6$&$48$&&$17$&$56$&
\end{tabular}
\endgroup
\end{center}

Für \textit{Ceres} haben wir nach den neuesten Elementen für 1803
\[
\begin{aligned}
e&=0.0788941.\\
k&=2.750681\\
g&=245^0\;34'\;56''
\end{aligned}
\]
\SourceLeaf{122}
Für die Erde gelten $e'$ und $k'$ wie oben sowie $g'=198^0\;35'\;31''$.
Daraus ergibt sich die Bedingungsgleichung
\[
\cos t'=.\,0.363444\cos t-0.004063
\]

Die Werte von $t'$ liegen also zwischen $68^0\;56'\;16''$ und
$111^0\;33'\;43''$ sowie zwischen $248^0\;26'\;17''$ und
$291^0\;3'\;44''$. Die heliozentrische Länge der Erde liegt folglich
zwischen $149^0\;55'$ und $192^0\;32'$ sowie zwischen $329^0\;25'$
und $12^0\;2'$. Die Erde kann \textit{Ceres} daher ungefähr nur
vom 19.\,Februar bis 3.\,April auf ihren nördlichen und vom
23.\,August bis 6.\,Oktober auf ihren südlichen Grenzen sehen.
Mit dieser Formel wurden wie zuvor bei \textit{Pallas} auf jeder
Grenze des Zodiakus von \textit{Ceres} 36 Punkte in Rektaszension
und Deklination berechnet. In Bezug auf den mehrfach erwähnten
Aufsatz wurden dabei für die Konstanten $a$, $A$ usw. die folgenden
Werte verwendet.
\[
\begin{array}{rcrrr}
a&=&79^0&30'&5''\\
b&=&114&42&11\\
c&=&27&7&24\\
A&=&170&49&4\\
B&=&85&42&29\\
C&=&59&36&33
\end{array}
\]
Bezeichnet $v$ die wahre Anomalie von \textit{Ceres}, so ergeben
sich die drei Koordinaten aus den folgenden Formeln.
\[
\begin{aligned}
x&=\frac{\alpha\sin(v+56^0\;24'\;0'')}{1-e\cos v}\\[8pt]
y&=\frac{\ObservedBeta\sin(v+331^0\;17'\;25'')}{1-e\cos v}\\[8pt]
z&=\frac{\gamma\sin(v+305^0\;11'\;29'')}{1-e\cos v}
\end{aligned}
\]
wobei
\[
\begin{aligned}
\log\alpha&=\log k\sin a=0.432108\\
\log\ObservedBeta&=\log k\sin b=0.397758\\
\log\gamma&=\log k\sin c=0.098319
\end{aligned}
\]
\SourceLeaf{123}

Statt dieser $72$ Punkte geben wir hier in der folgenden Tabelle
nur die $144$ Punkte an, die auf ähnliche Weise wie bei
\textit{Pallas} zwischen ihnen eingefügt wurden.

\begin{center}
\begingroup\small\setlength{\tabcolsep}{3pt}
\begin{tabular}{r|rrl|rrl||r|rrl|rrl}
\multicolumn{14}{c}{\textit{Zodiakus der Ceres}}\\[4pt]
&\multicolumn{6}{c||}{Deklination}&&\multicolumn{6}{c}{Deklination}\\
\shortstack{Rektasz.}&\multicolumn{3}{c|}{\shortstack{der nördlichen\\Grenze}}&\multicolumn{3}{c||}{\shortstack{der südlichen\\Grenze}}&\shortstack{Rektasz.}&\multicolumn{3}{c|}{\shortstack{der nördlichen\\Grenze}}&\multicolumn{3}{c}{\shortstack{der südlichen\\Grenze}}\\\hline
$0^0$&$8^0$&$37'$&südl.&$17^0$&$29'$&südl.&$180^0$&$18^0$&$49'$&nördl.&$8^0$&$17'$&nördl.\\
$5$&$6$&$15$&&$15$&$13$&&$185$&$16$&$31$&&$5$&$56$&\\
$10$&$3$&$45$&&$12$&$51$&&$190$&$14$&$5$&&$3$&$28$&\\
$15$&$1$&$11$&&$10$&$26$&&$195$&$11$&$33$&&$0$&$55$&\\
$20$&$1$&$27$&nördl.&$7$&$56$&&$200$&$8$&$57$&&$1$&$41$&südl.\\
$25$&$4$&$7$&&$5$&$18$&&$205$&$6$&$16$&&$4$&$19$&\\
$30$&$6$&$48$&&$2$&$38$&&$210$&$3$&$34$&&$6$&$56$&\\
$35$&$9$&$26$&&$0$&$3$&&$215$&$0$&$52$&&$9$&$32$&\\
$40$&$12$&$1$&&$2$&$30$&nördl.&$220$&$1$&$47$&südl.&$12$&$3$&\\
$45$&$14$&$31$&&$4$&$58$&&$225$&$4$&$21$&&$14$&$29$&\\
$50$&$16$&$54$&&$7$&$20$&&$230$&$6$&$50$&&$16$&$48$&\\
$55$&$19$&$9$&&$9$&$35$&&$235$&$9$&$12$&&$18$&$59$&\\
$60$&$21$&$16$&&$11$&$42$&&$240$&$11$&$25$&&$21$&$1$&\\[3pt]
$65$&$23$&$12$&&$13$&$39$&&$245$&$13$&$28$&&$22$&$53$&\\
$70$&$24$&$58$&&$15$&$27$&&$250$&$15$&$21$&&$24$&$35$&\\
$75$&$26$&$34$&&$17$&$4$&&$255$&$17$&$3$&&$26$&$6$&\\
$80$&$27$&$59$&&$18$&$30$&&$260$&$18$&$34$&&$27$&$26$&\\
$85$&$29$&$13$&&$19$&$45$&&$265$&$19$&$53$&&$28$&$36$&\\
$90$&$30$&$16$&&$20$&$49$&&$270$&$21$&$1$&&$29$&$35$&\\[3pt]
$95$&$31$&$8$&&$21$&$42$&&$275$&$21$&$56$&&$30$&$24$&\\
$100$&$31$&$50$&&$22$&$24$&&$280$&$22$&$40$&&$31$&$1$&\\
$105$&$32$&$21$&&$22$&$55$&&$285$&$23$&$13$&&$31$&$29$&\\
$110$&$32$&$42$&&$23$&$14$&&$290$&$23$&$34$&&$31$&$46$&\\
$115$&$32$&$51$&&$23$&$22$&&$295$&$23$&$43$&&$31$&$52$&\\
$120$&$32$&$50$&&$23$&$19$&&$300$&$23$&$41$&&$31$&$48$&\\[3pt]
$125$&$32$&$40$&&$23$&$5$&&$305$&$23$&$28$&&$31$&$34$&\\
$130$&$32$&$18$&&$22$&$40$&&$310$&$23$&$4$&&$31$&$9$&\\
$135$&$31$&$46$&&$22$&$3$&&$315$&$22$&$27$&&$30$&$34$&\\
$140$&$31$&$3$&&$21$&$16$&&$320$&$21$&$40$&&$29$&$49$&\\
$145$&$30$&$9$&&$20$&$17$&&$325$&$20$&$40$&&$28$&$53$&\\
$150$&$29$&$4$&&$19$&$6$&&$330$&$19$&$30$&&$27$&$46$&\\[3pt]
$155$&$27$&$49$&&$17$&$45$&&$335$&$18$&$8$&&$26$&$29$&\\
$160$&$26$&$22$&&$16$&$12$&&$340$&$16$&$35$&&$25$&$1$&\\
$165$&$24$&$44$&&$14$&$28$&&$345$&$14$&$51$&&$23$&$23$&\\
$170$&$22$&$56$&&$12$&$34$&&$350$&$12$&$55$&&$21$&$34$&\\
$175$&$20$&$57$&&$10$&$30$&&$355$&$10$&$51$&&$19$&$36$&\\
$180$&$18$&$49$&&$8$&$17$&&$360$&$8$&$37$&&$17$&$29$&\\
\end{tabular}
\endgroup
\end{center}
